\section*{Abstract}
\label{sec:abstract}
\addcontentsline{toc}{section}{\nameref{sec:abstract}}

Markdown's syntax is small and easily readable, which makes it well suited to writing documentation: being plain text, it can be opened and edited in a regular text editor, and its changes can easily be tracked by version control systems.
Using it for writing teaching materials is not ideal, though, not because Markdown itself falls short, but because the Moodle learning management system does not support it.
In Moodle, material can be authored as H5P packages in a browser-based editor.
This, however, means content written in Markdown is ported by hand, and the port is repeated every time the source changes.
H5P packages can be created for different content types, the most apt of which for representing Markdown files is Interactive Book: it represents content as chapters, each chapter made up of content blocks.
Those blocks come from a small set of content types with fixed nesting rules, so some Markdown constructs -- an image inside a table cell, for instance -- have no direct counterpart.
This thesis derives the limits of H5P's Interactive Book type and presents md2h5p, a converter that turns a folder of Markdown files into an H5P Interactive Book.
It renders each Markdown file into a chapter, finds alternatives for the Markdown constructs H5P does not support directly, and warns about constructs that can only be represented in a degraded form or not at all.
An extension adds single- and multiple-choice questions, written as Markdown fenced code blocks, and a
merge step folds a fresh conversion back into a package that has since been edited in the H5P editor.
Teachers can therefore keep authoring in Markdown, revise it after publication, and still
hand their students content the platform can embed.
